% 钱德拉塞卡拉·拉曼（综述）
% license CCBYSA3
% type Wiki

本文根据 CC-BY-SA 协议转载翻译自维基百科\href{https://en.wikipedia.org/wiki/C._V._Raman}{相关文章}

钱德拉塞卡拉·文卡塔·“C. V.”·拉曼爵士（Sir Chandrasekhara Venkata "C. V." Raman，/ˈrɑːmən/，\(^\text{[2]}\)音译“拉曼”；泰米尔语：சந்திரசேகர வெங்கட ராமன்，罗马化：Cantiracēkara Veṅkaṭa Rāmaṉ；1888年11月7日—1970年11月21日）是一位印度物理学家\(^\text{[3]}\)，以其在光散射领域的研究而闻名。\(^\text{[4]}\)他使用自己研制的分光仪，与学生K. S. 克里希南共同发现：当光线穿过透明物质时，被偏转的光会改变其波长。这种此前未知的光散射现象，他们称之为“修饰散射”，后来被命名为“拉曼效应”或“拉曼散射”。1930年，拉曼因这一发现获得诺贝尔物理学奖，成为第一位获得该奖项的亚洲人和非白人。\(^\text{[5]}\)
